\documentclass[11pt]{article}
\usepackage[utf8]{inputenc}
\usepackage{amsmath,amssymb,amsthm}
\usepackage[margin=1in]{geometry}
\usepackage{hyperref}
\usepackage{xcolor}
\usepackage{newunicodechar}
\newunicodechar{ℝ}{\ensuremath{\mathbb{R}}}
\newunicodechar{⟨}{\ensuremath{\langle}}
\newunicodechar{⟩}{\ensuremath{\rangle}}
\newunicodechar{Γ}{\ensuremath{\Gamma}}
\newunicodechar{·}{\ensuremath{\cdot}}
\newunicodechar{²}{\ensuremath{^2}}
\newunicodechar{³}{\ensuremath{^3}}
\newunicodechar{⁴}{\ensuremath{^4}}
\newunicodechar{⁶}{\ensuremath{^6}}
\newunicodechar{κ}{\ensuremath{\kappa}}
\newunicodechar{≤}{\ensuremath{\leq}}
\newunicodechar{−}{\ensuremath{-}}
\newunicodechar{∼}{\ensuremath{\sim}}
\newunicodechar{μ}{\ensuremath{\mu}}
\title{The sixth cumulant of the even-monomial Gibbs weight}
\author{Tide retrospective, laplace seabed}
\date{2026-09-06}
\begin{document}
\sloppy
\maketitle
\section{Setting}
The laplace monomial track carries a cumulant ladder against the symmetric weight
$\exp(-t\,x^{2k}/(2k)!)$: the variance $\kappa_2$ and the excess kurtosis $\kappa_4$, each in Gamma
closed form and with its $t^{-\ldots}$ asymptotic scaling. This tide adds the next even rung, the
sixth cumulant $\kappa_6$, completing the ladder at the clean frontier (the general even cumulant
would require set-partition / Bell-polynomial machinery).
\section{The thing formalised}
For a symmetric weight the mean and all odd moments vanish, so
$\kappa_6 = \langle x^6\rangle - 15\langle x^4\rangle\langle x^2\rangle + 30\langle x^2\rangle^3$, and
\[
  \kappa_6 = \Big(\tfrac{(2k)!}{t}\Big)^{3/k}\big(r_3 - 15\,r_2 r_1 + 30\,r_1^3\big),\qquad
  r_j := \frac{\Gamma((2j+1)/(2k))}{\Gamma(1/(2k))},
\]
hence $\kappa_6 \sim K_6(k)\,t^{-3/k}$ at $\mathrm{atTop}$. The Lean file gives the closed form
\texttt{monomial\_sixth\_cumulant}, the constant \texttt{sixthCumulantConst}, and the three asymptotic
rungs. At $k=1$ the bracket is $15 - 15\cdot 3 + 30 = 0$: a Gaussian has vanishing sixth cumulant.
\section{Proof strategy}
The closed form is the three-moment mirror of the kurtosis proof. Rewrite $x^6, x^4, x^2$ into the
even form $x^{2\cdot 3}, x^{2\cdot 2}, x^{2\cdot 1}$ (definitional), substitute the even-moment closed
form at $j=3,2,1$, and reduce everything to the common base $B := ((2k)!/t)^{1/k}$. Two power folds do
this: $((2k)!/t)^{2/k} = B^2$ (as in the kurtosis) and $((2k)!/t)^{3/k} = B^3$. The cube is the one
new step: \texttt{Real.rpow\_natCast} turns the natural power $B^3$ into the real power $B^{(3:\mathbb{R})}$,
and \texttt{Real.rpow\_mul} collapses $(base^{1/k})^{3} = base^{(1/k)\cdot 3}$; a \texttt{push\_cast}
and \texttt{ring} match the exponent. With both folds the goal is polynomial in $B$ and the Gamma
ratios (opaque atoms), so \texttt{ring} closes it. The three asymptotic rungs are the constant-sequence
and reflexivity arguments reused verbatim from the moment and kurtosis asymptotics.
\section{Roadblocks and resolutions}
None; the tide compiled on the first attempt. The only forethought was including an \texttt{Authors:}
line in the header from the outset, having learned the previous tide that the \texttt{style.header}
linter flags its absence.
\section{What was Mathlib, what was new}
Mathlib supplied the \texttt{Real.rpow} algebra (including \texttt{rpow\_natCast}/\texttt{rpow\_mul} for
the cube) and the \texttt{IsEquivalent}/\texttt{Tendsto} API. The closed-form even moments were already
in the seabed. New here is the sixth cumulant, its constant, and its scaling --- the last clean rung of
the monomial cumulant ladder. No GPT consult was needed.
\end{document}
